\documentclass[11pt]{article}
\usepackage{amsmath,amssymb}
\title{Exact affine Gaussian chain: arbitrary forcing and retained holonomy}\author{}\date{}
\begin{document}\maketitle
Let $O\in SO(3)$ and consider actual Lebesgue kernels
\[
 K_t(x,y)=\exp[-\tfrac32(|x|^2+|y|^2)+x\cdot y+b_t\cdot x+c_t\cdot y+d_t].
\]
There is no smallness condition on these finite vectors. The period closes with
$x_m=Ox_0$, and $C_Of(x)=f(O^{-1}x)$. Thus its trace is the path integral of
$W=K_0\cdots K_{m-1}C_O$. Endpoint coefficients are interpreted literally with
this twisted identification; covariant repetition gives the path integral of
$W^q$, not an independent untwisted replication of the coefficients.

The quadratic path action has Hessian $L=6I-S-S^*$, where $S$ is the orthogonal
twisted cyclic shift. Hence $L\ge4I$. Its unique minimizer $z$ solves
\[
 6z_t-z_{t-1}-z_{t+1}=b_t+c_{t-1},\qquad z_m=Oz_0,
 \qquad \max_t|z_t|\le\tfrac14\max_t|b_t+c_{t-1}|.
\]
The last estimate follows at a maximum norm index, by taking its inner product
with $z_t/|z_t|$; it also follows from the Neumann series for $L^{-1}$.
All seam indices carry the appropriate orthogonal transport.

Translate $x_t=z_t+y_t$. Stationarity cancels every linear term, including the
seam. With
\[
 F_m=\sum_t d_t+\tfrac12\langle b+c_{-1},L^{-1}(b+c_{-1})\rangle,
\]
the exact period trace is
\[
 \operatorname{Tr}W=e^{F_m}\operatorname{Tr}(H^m C_O),\quad
 H(x,y)=e^{-\frac32(|x|^2+|y|^2)+x\cdot y}.
\]
The minimizer of $q$ repeated periods is the covariant repetition of $z$:
it is a stationary point, and strict convexity proves it is the unique minimum.
The scalar is therefore EXACTLY $qF_m$, not a sum of approximate local vacua.
Consequently
\[
 \operatorname{Tr}W^q=e^{qF_m}\operatorname{Tr}(H^{mq} C_{O^q}).                 \tag{1}
\]
No error depends on the size of the affine forcing.

For $\ell_0=[\pi/(3/2+\sqrt2)]^{3/2}$ and $r_0=3-\sqrt8$, the Mehler spectrum
on homogeneous oscillator excitations is $\ell_0r_0^n$ with the natural
$SO(3)$ symmetric-power action. Summing its generating function gives
\[
 \operatorname{Tr}(H^L C_R)=\frac{\ell_0^L}{\det(I-r_0^L R)},\qquad R\in SO(3).
                                                                            \tag{2}
\]
This identity can alternatively be obtained by completing the finite cyclic
Gaussian integral; it retains all three oscillator coordinates and the literal
endpoint rotation. In particular its denominator is positive and tends to one
uniformly in $R$. Taking $q$th roots in(1), or using the positive compact kernel's
Perron asymptotic, identifies the EXACT period eigenvalue
\[
 \rho(W)=e^{F_m}\ell_0^m,\qquad
 \frac{\operatorname{Tr}W^q}{\rho(W)^q}
       =\det(I-r_0^{mq}O^q)^{-1}.                                          \tag{3}
\]
An equivalent formal kernel conjugation uses $x\mapsto x-z_t$ and exponential
linear multipliers whose neighboring factors telescope. These multipliers can
be unbounded on full Euclidean $L^2$; (1)--(3) were instead justified by finite
Gaussian integrals and their exact traces, so no bounded-similarity claim is
needed. On a compact group the corresponding finite-cutoff multipliers are
bounded, though not uniformly in the cutoff.

This calculation identifies what can cancel for a path-centered compact
comparison: the classical path scalar and the linear forces. It does not by
itself prove a Wilson estimate or a probability bound for any bath class.
\end{document}
